\scanpage{173}
\mspage{122}
% archive sheet number: 173
\unsure{Il} faut entendre par là que (pour $W$ \inserted{$\subset\mathrm{Fl}(\mathrm{Cat})$} donné) la classe \inserted{$\mathrm{hot}_W(\unsure{\mathscr{X}|\mathcal X})$} de $\mathscr{X}$ dans $\mathrm{Hot}_W$ a un \lq\lq sens homotopique\rq\rq{} remarquable \inserted{en termes de ($*$)} ; de façon précise, le diagramme ($*$) définit un objet $\unsure{\xi|\zeta}$ de $\mathrm{HOT}_W(\unsure{\mathscr{I}|\Phi})$ (qui est une catégorie de fractions de la catégorie de tous les diagrammes du type ($*$)\ldots), et si \deleted{\unsure{$\xi$}}

\[
\pi:\mathscr{I}\longrightarrow e
\]
désigne la projection absolue de $\mathscr{I}$, on a

\[
\pi_!(\xi)=\deleted{\unsure{h|H}}\mathrm{hot}_W(\mathscr{X}),
\]
d'après le calcul connu du foncteur $\pi_!$, adjoint à gauche du foncteur $\pi^*$

\[
\mathrm{HOT}_W(e)\overset{\text{déf}}{=}\mathrm{Hot}_W\longrightarrow\mathrm{HOT}_W(\mathscr{I}).
\]

% slanted
\margin{Cela suppose qu'on a $W(\unsure{3|2})$, non \unsure{simplem\textsuperscript{t}|seulement} $W(\unsure{2|3}\,\mathrm{bis})$.}

Or pour tout objet $I$ de $\mathrm{Cat}$, de projection absolue $p_I:I\to e$, le foncteur

\[
p_{I\,!}:\mathrm{HOT}_W(I)\unsure{-|\to}\mathrm{HOT}_W(e)=\mathrm{Hot}_W
\]
doit être regardé comme l'équivalent \lq\lq homotopique\rq\rq{} \unsure{,} du foncteur limite inductive $\underset{I}{\varinjlim}$ \unsure{des|de} diagrammes de type $I^\circ$ dans la catégorie $\mathrm{Hot}_W$ -- lesquelles limites inductives \uline{n'existent pas} dans $\mathrm{Hot}_W$ (sauf
